\section{SPT Phase Labels and Virtual-Boundary Nondecay}
%% ---------------------------------------------------------

For one-dimensional gapped systems within the injective symmetric MPS and
parent-Hamiltonian framework~\cite{Schuch2011Phases}, symmetry-protected
topological (SPT) phases are classified by the second cohomology group
$H^2(G,\mathrm{U}(1))$ of the symmetry group, appearing as the cohomology class
of the virtual projective representation.

\begin{definition}[Equality of virtual cocycle labels]
    \label{def:is_same_spt_phase}
    \lean{MPSTensor.IsSameSPTPhase}
    \leanok
    \uses{def:cohomologous_to, def:projective_rep}
    For tensors $A,B$ and a fixed physical representation
    $U:G\to\GL_d(\C)$, the predicate called \emph{same SPT phase}
    asserts the existence of factor systems
    $\omega_A,\omega_B:G\times G\to\C^\times$ and projective bond
    representations $\rho_A,\rho_B$ such that
    \begin{align}
        \widetilde A_g^i
         & =\rho_A(g^{-1})A^i\rho_A(g^{-1})^{-1}, \\
        \widetilde B_g^i
         & =\rho_B(g^{-1})B^i\rho_B(g^{-1})^{-1},
        \notag
    \end{align}
    and $\omega_A\sim\omega_B$.
\end{definition}

\begin{remark}[Cohomology with $\C^\times$ and $\mathrm{U}(1)$ coefficients]
    \label{rem:spt_scalar_character_conventions}
    \uses{def:is_same_spt_phase, def:h2_cx,
        thm:mpug_circle_complex_units_h2}
    While factor systems here take values in $\C^\times$, for finite groups $G$
    every class admits a unit-modulus representative, identifying
    $H^2(G,\C^\times)$ with $H^2(G,\mathrm U(1))$
    (Theorem~\ref{thm:mpug_circle_complex_units_h2}).
\end{remark}

\begin{theorem}[Classification along symmetric MPS parent-Hamiltonian paths]
    \label{thm:spt_classification}
    \notready
    \uses{def:is_same_spt_phase, def:h2_cx}
    For $p=0,1$, let $A_p$ be an injective MPS tensor whose parent Hamiltonian
    $H_p$ has a unique ground state and is symmetric under a unitary
    representation $U^p$ of $G$ on $\mc H_p$. Choose the endpoint phase
    gauges. After a common finite blocking, the embedded endpoints admit a
    continuous symmetric MPS path with corresponding translationally invariant,
    finite-range parent Hamiltonians that vary smoothly and remain uniformly
    gapped if and only if their virtual projective representations determine
    the same class in $H^2(G,\mathrm U(1))$.

    More precisely, for the forward construction there exist a path-sector
    Hilbert space $\mc H_{\mathrm{path}}$ and a linear representation
    $U^{\mathrm{path}}$ such that the endpoints embed into
    \begin{align}
        \mc H & =\mc H_0\oplus\mc H_1\oplus\mc H_{\mathrm{path}}, &
        U_g   & =U_g^0\oplus U_g^1\oplus U_g^{\mathrm{path}}.
        \notag
    \end{align}
    The converse quantifies over smooth paths of MPS with corresponding
    parent Hamiltonians \cite[Section~II.F, lines~929--953]{Schuch2011Phases}.
\end{theorem}

\begin{proof}
    First deform each blocked injective MPS to its isometric form; the
    deformation preserves the symmetry relation
    \begin{align}
        U_g\mc P & =\mc P(V_g\otimes\overline{V_g}).
        \notag
    \end{align}
    Along any continuous symmetric MPS path quantified in the theorem, the
    tensor and its induced virtual action $V_g$ vary continuously. The
    cohomology class is discrete, so it cannot change along that path. This
    proves necessity within the MPS/parent-Hamiltonian framework.

    Conversely, let the endpoint isometric forms have virtual actions
    $V_g^0$ and $V_g^1$ in the same cohomology class.  Embed their bond spaces
    into the direct sum and interpolate between the maximally entangled bonds
    by
    % Local fix: the second range ends at D_0+D_1, not D_1; see
    % docs/paper-gaps/spc11_spt_interpolation_upper_range.tex.
    \begin{align}
        |\omega(\gamma)\rangle
         & =(1-\gamma)\sum_{i=1}^{D_0}|i,i\rangle
        +\gamma\sum_{i=D_0+1}^{D_0+D_1}|i,i\rangle.
        \notag
    \end{align}
    The commuting parent Hamiltonians along this interpolation remain gapped.
    Their symmetry is
    \begin{align}
        (V_g^0\oplus V_g^1)\otimes
        (\overline{V_g^0}\oplus\overline{V_g^1})
         & =\widehat U_g^0\oplus\widehat U_g^1
        \oplus\widehat U_g^{\mathrm{path}},    \\
        \widehat U_g^{\mathrm{path}}
         & =(V_g^0\otimes\overline{V_g^1})
        \oplus(V_g^1\otimes\overline{V_g^0}).
        \notag
    \end{align}
    Equality of the cohomology classes is exactly what makes the cross-term
    action $\widehat U^{\mathrm{path}}$ a linear representation.  Conjugating
    back by the endpoint isometries gives the common direct-sum physical
    representation and the required symmetric gapped path.
    % The Hamiltonian-path and isometric-form bridges are not yet represented
    % by Lean declarations; the \notready marker records that status.
\end{proof}

\begin{theorem}[Virtual-boundary nondecay for injective MPV-symmetric tensors]%
    \label{thm:has_string_order_of_symmetric_injective}
    \lean{MPSTensor.hasStringOrder_of_symmetric_injective}
    \leanok
    \uses{def:has_string_order, thm:virtual_rep_injective}
    Let $A$ be an injective tensor, symmetric under a unitary
    representation $U$. Assume $\E_A(\Id)=\Id$, $\Lambda$ is positive
    definite, $\tr(\Lambda)=1$, and $\E_A^\dagger(\Lambda)=\Lambda$.
    Then the virtual-boundary nondecay predicate holds for every group
    element $g$.
\end{theorem}

\begin{proof}\leanok
    \uses{thm:virtual_rep_injective, thm:twisted_transfer_modulus_one_gauge_phase, thm:virtual_unitary_of_twisted_gauge_phase, thm:boundary_state_invariant_of_virtual_unitary, lem:has_string_order_of_local_symmetry}
    The virtual representation theorem produces
    $\rho(g^{-1}) \in \GL_{D}(\C)$ satisfying
    $\sum_j U(g)_{ij}\,A^j
        = \rho(g^{-1})\,A^i\,\rho(g^{-1})^{-1}$.
    A direct computation using $\E_A(\Id)=\Id$ shows that
    $\E_{U(g)}(\rho(g^{-1}))=\rho(g^{-1})$, so
    $\rho(g^{-1})$ is a nonzero eigenvector with eigenvalue~$1$.
    Theorem~\ref{thm:twisted_transfer_modulus_one_gauge_phase} gives
    gauge-phase equivalence with the companion tensor, and
    Theorem~\ref{thm:virtual_unitary_of_twisted_gauge_phase} extracts a virtual
    unitary and phase.
    By Theorem~\ref{thm:boundary_state_invariant_of_virtual_unitary}, this
    unitary preserves the boundary state, so
    Lemma~\ref{lem:has_string_order_of_local_symmetry} gives
    virtual-boundary nondecay.
\end{proof}

\begin{remark}[Virtual-boundary nondecay does not determine the SPT label]%
    \label{rem:string_order_symmetry_diagnostic}
    \uses{thm:has_string_order_iff_of_symmetric_injective}
    Theorem~\ref{thm:has_string_order_of_symmetric_injective} applies to
    every injective symmetric tensor with canonical normalisations,
    irrespective of its cocycle class, so tensors in different SPT phases
    all satisfy virtual-boundary nondecay for every group element.  This
    nondecay is equivalent to virtual local covariance
    (Theorem~\ref{thm:string_order_iff_local_symmetry}); the invariant
    that separates SPT phases is the cohomology class of the virtual
    representation cocycle (Theorem~\ref{thm:spt_classification}).
\end{remark}

\subsection{Degeneracy of the entanglement spectrum}
\label{sec:spt_entanglement_spectrum}

For an injective tensor in canonical form, the positive fixed point of the
transfer map is the half-chain reduced density matrix, whose spectrum is the
entanglement spectrum.  A symmetry forces this fixed point to commute with the
virtual gauges, and a non-commuting projective representation then forces every
eigenvalue to be degenerate \cite[Section~III.A]{Cirac2021Matrix}.  Here the
canonical form is $\E_A(\Id)=\Id$ and the fixed point is that of the adjoint
map $\E_A^\dagger$; the reflected normalization, in which the fixed point of
$\E_A$ itself is used, describes the same spectrum.

\begin{lemma}[Anticommuting invertible operators force even dimension]
    \label{lem:even_finrank_of_anticommute}
    \lean{Module.End.even_finrank_of_anticommute, Matrix.even_finrank_eigenspace_of_anticommute}
    \leanok
    Let $f,g$ be invertible operators on a finite-dimensional complex vector
    space $V$ with $fg=-gf$.  Then $\dim V$ is even.  Consequently, if a matrix
    $\Lambda$ commutes with two anticommuting invertible matrices $X$ and $Y$,
    every eigenspace of $\Lambda$ has even dimension.
\end{lemma}

\begin{proof}\leanok
    Taking determinants in $fg=-gf$ gives
    $\det f\det g=(-1)^{\dim V}\det g\det f$ with $\det f\det g\neq0$, so
    $(-1)^{\dim V}=1$.  For the second statement, $X$ and $Y$ map each eigenspace
    of $\Lambda$ into itself, and their restrictions are invertible and
    anticommute.
\end{proof}

\begin{theorem}[Virtual gauges fix the boundary state]
    \label{thm:virtual_gauge_fixes_boundary_state}
    \lean{MPSTensor.exists_unitary_gauge_of_symmetry, MPSTensor.commute_boundaryState_of_symmetry}
    \leanok
    \uses{def:injective, lem:gauge_phase_unique, thm:virtual_unitary_of_twisted_gauge_phase, thm:boundary_state_invariant_of_virtual_unitary}
    Let $A$ be injective with $\E_A(\Id)=\Id$, let $\Lambda$ be a positive
    definite matrix with $\E_A^\dagger(\Lambda)=\Lambda$, and let $u$ be a
    unitary with
    $\sum_j u_{ij}A^j=\zeta\,XA^iX^{-1}$ for all $i$, where $X\in\GL_{D}(\C)$
    and $\zeta\neq0$.  Then $X=c\,W$ for a nonzero scalar $c$ and a unitary $W$
    with $W\Lambda W^\dagger=\Lambda$.  In particular $X\Lambda=\Lambda X$.
\end{theorem}

\begin{proof}\leanok
    \uses{lem:gauge_phase_unique, thm:virtual_unitary_of_twisted_gauge_phase, thm:boundary_state_invariant_of_virtual_unitary}
    Theorem~\ref{thm:virtual_unitary_of_twisted_gauge_phase}, applied to the
    physical unitary $u^\dagger$, gives a unitary $V$ and a phase $\mu$ with
    $\sum_j (u^\dagger)_{ij}A^j=\mu\,VA^iV^\dagger$, and
    Theorem~\ref{thm:boundary_state_invariant_of_virtual_unitary} gives
    $V^\dagger\Lambda V=\Lambda$ after normalizing the trace of $\Lambda$.
    Undoing the physical mixing yields $A^i=\mu\zeta\,(VX)A^i(VX)^{-1}$, so
    Lemma~\ref{lem:gauge_phase_unique} makes $VX$ a nonzero scalar.  Hence $X$
    is a multiple of $W=V^\dagger$.
\end{proof}

\begin{theorem}[Even degeneracy of the entanglement spectrum]
    \label{thm:spt_entanglement_spectrum_even}
    \lean{MPSTensor.even_finrank_eigenspace_of_anticommuting_gauges}
    \leanok
    \uses{thm:virtual_gauge_fixes_boundary_state, lem:even_finrank_of_anticommute}
    In the setting of Theorem~\ref{thm:virtual_gauge_fixes_boundary_state}, let
    $u,v$ be unitaries with virtual gauges $X,Y$, that is,
    $\sum_j u_{ij}A^j=\zeta\,XA^iX^{-1}$ and
    $\sum_j v_{ij}A^j=\eta\,YA^iY^{-1}$ with $\zeta,\eta\neq0$.  If $XY=-YX$,
    then every eigenspace of $\Lambda$ has even dimension: every eigenvalue of
    the entanglement spectrum is degenerate \cite[Section~III.A]{Cirac2021Matrix}.
\end{theorem}

\begin{proof}\leanok
    \uses{thm:virtual_gauge_fixes_boundary_state, lem:even_finrank_of_anticommute}
    Both gauges commute with $\Lambda$ by
    Theorem~\ref{thm:virtual_gauge_fixes_boundary_state}; apply
    Lemma~\ref{lem:even_finrank_of_anticommute}.
\end{proof}
